\section{Lesson 7, 06/10}
\subsection{Classical Field theory}
Before delving into QFT, it is necessary to understand classical field theory.

A classical system is described by the generalized coordinates
$$q_i(t),\;\text{with }i=1,2...N$$
where $N$ is the (finite) number of degrees of freedom and $t$ is a parameter.
From this the Lagrangian is constructed
$$L(q,\dot q).$$
A field theory replaces the set of the coordinates$\left\{q_i\right\}$ with a
\textit{continuous function of space-time}, called a \textbf{field}:
$$\phi(t,\mathbf{x}),\;\mathbf{x}\in\mathbb{R}^3$$
the continuous coordinates $\mathbf{x}$ replace the discrete index $i$ and so the
degrees of freedom become infinite.

The next step is to construct a Lagrangian density
$$\mathcal{L}(\phi,\pd_\mu\phi)$$
and solve for the equations of motion.

\textbf{Example:} from finite to infinite degrees of freedom.\\
Suppose a system of $N$ point-like masses interacting with an elastic potential:
the Lagrangian will be
$$L=\sum_{k=1}^{N}\left[\frac{1}{2}m_i\dot x_i^2-\frac{1}{2}k(x_i-x_{i+1})^2\right].$$
Consider $m_i = m$ and the distance between them to be constant and equal to
$\varepsilon$:
$$L = \sum_{k=1}^{N}\frac{\varepsilon}{2}\left[\frac{m}{\varepsilon}\dot x_i^2
    -k\frac{(x_i-x_{i+i})^2}{\varepsilon}\right]$$
The limit as $\varepsilon$ goes to 0 is
$$\lim\limits_{\varepsilon\to0} L = \int_{}^{}\dd x \frac{1}{2}\left[
\mu\dot\phi(x)^2 - Y\left(\pdv[]{\phi}{x}\right)^2\right]$$
where $\frac{m}{\varepsilon}\rightarrow\mu$, $k\varepsilon\rightarrow Y$
(which is a Young module) and $x_i(t)\rightarrow\phi(t,x)$.
Applying the least action principle yields the Euler-Lagrange equation
$$\pdv[2]{\phi}{x} - \frac{\mu}{Y}\pdv[2]{\phi}{t} = 0.$$\\

In general, the Lagrangian density is a function of the spacetime coordinates $x^\mu$
through the dependency on a field and its derivatives
$$\mathcal{L}(x^\mu) = \mathcal{L}(\phi,\pd_mu\phi)$$
and the Lagrangian is its integral on the physical space
$$L(t) = \int_{}^{}\dd^3 x \mathcal{L}(\phi,\pd_\mu\phi);$$
another important quantity is the action, i.e. the integral over spacetime of the
density:
$$S = \int_{}^{}\dd^4x \mathcal{L}(\phi,\pd_\mu\phi).$$\\

All the theories constructed must be \textit{local}, i.e. they must obey causality,
so in the Lagrangian density $\mathcal{L}$ will never present terms like
$$\phi(x^\mu)\phi(y^\mu),\;\text{with }x^\mu\neq y^\mu$$
which represents an instantaneous interaction between two points of spacetime.

\textbf{Def.} A field theory is local when the dependency con $x^\mu$ is
$$\mathcal{L}(x^\mu) =\mathcal{L}(\phi(x^\mu),\pd_\mu\phi(x^\mu))$$
which means that the dynamics of the system do not connect different points in space
instantaneously.

\subsubsection{Transformations of fields}
Since the theory must be covariant, it is pivotal to understand how fields
behave under Poincaré transformations: consider a point in spacetime expressed in
coordinates as $x^\mu$ and the value of a field at that point $f(x^\mu)$; consider
then an infinitesimal transformation $x'^\mu = x^\mu+\delta x^\mu$ so that
$$\Delta f = f'(x'^\mu) - f(x^\mu) = f'(x^\mu+\delta x^\mu) - f(x^\mu) = f'(x^\mu)
-f(x^\mu) + \pd_\mu f'(x^\mu)\delta x^\mu + O(\delta^2 x^\mu)
$$
at first order, noting that $\pd_\mu f'(x^\mu) = \pd_\mu f(x^\mu) + O(\delta x^\mu)$,
this can be written as
$$\Delta f = \delta f +\pd_\mu f(x^\mu)\delta x^\mu.$$
This can be seen as the definition of the operator $\Delta$ on fields
$$\Delta = \delta + \delta x^\mu \pd_\mu:$$
under a Poincaré transformation, the variation of a field is made up of the
functional variation in the field itself ($\delta$) and the variation in the field
given by the transformation of the coordinates (called transport term).

\textbf{Example:} spacetime translations.\\
Consider the infinitesimal Poincaré transformation
$$x^\mu\rightarrow x'^\mu = x^\mu+\varepsilon^\mu$$
by definition, a local field $f$ is invariant under translations, so
$$\Delta f = 0$$
and this means
$$\delta f = - \varepsilon^\mu\pd_\mu f$$
and since $P_\mu = i\pd_\mu$, 
$$\delta f =-i\varepsilon^\mu P_\mu f.$$
This means that this functional transformation of the field is generated by the
4-momenta
$$f' = \exp\left[-i\varepsilon^\mu P_\mu\right]f$$

The Lorentz transformations, on the other hand, depend on the representation the field
belongs to: the simplest is the scalar representation.

\textbf{Example:} Lorentz transformation of a scalar field.\\
A scalar field is a field invariant under Poincaré transformations
$$\phi(x^\mu) = \phi'(x'^\mu).$$
As seen before, the translations are easy
$$\phi(x^\mu)\rightarrow\phi'(x'^\mu) = \exp\left[-i\varepsilon^\mu P_\mu\right]
\phi(x^\mu).$$
Consider now a generic Lorentz transformation
$$x^\mu\rightarrow x'^\mu = x^\mu + \varepsilon^{\mu\rho}x_\rho$$
where $\varepsilon^{\mu\rho}$ are the 6 completely antisymmetric parameters of the
transformation.\\
By definition,
$$\Delta \phi = 0 = \delta\phi + \varepsilon^{\mu\rho}x_\rho\pd_\mu \phi$$
and so
$$\delta\phi = -\varepsilon^{\mu\rho}x_\rho\pd_\mu\phi \rightarrow \delta
=-\varepsilon^{\mu\rho}x_\rho\pd_\mu$$
this operator is non-zero only for the antisymmetric part of $x_\rho\pd_\mu$, so
antisymmetrizing yields
$$\delta = -\varepsilon^{\mu\rho}\frac{1}{2}\left(x_\rho\pd_\mu-x_\mu\pd_\rho\right)=
-\varepsilon^{\mu\rho}\frac{i}{2}L_{\rho\mu}
$$
where $L_{\rho\mu}$ is the generalized angular momentum expressed in its
infinite-dimensional form, related to the $M_{\rho\mu}$ generators.\\
\textbf{Observation:} no spin operator is found, as expected, because the scalar
field belongs to the $(0,0)$ representation, and so the only possible angular momentum
of the system is orbital.

And so the general Lorentz transformation on a scalar field is
$$\phi'(x'^\mu) = \exp\left[-\frac{i}{2}\varepsilon^{\mu\rho}L_{\rho\mu}\right]
\phi(x^\mu)$$

\textbf{Observation:} Any power of a scalar field is a Lorentz scalar (and so all
functions of $\phi$ expressed as power series), the derivative of $\phi$ is
not scalar, while the double derivative is scalar.\\

\textbf{Example:} Lorentz transformation of a vector field.\\
Consider the gradient vector field
$$\pd_\mu \phi,$$
where $\phi$ is a scalar field: what does the operator $\Delta$ do to this kind of
field?
$$\Delta\pd_\mu\phi = \delta\pd_\mu\phi+\delta x^\nu\pd_\nu\pd_\mu\phi$$
since $\Delta \phi =0$, this is also
$$\Delta\pd_\mu\phi = \left[\Delta,\pd_\mu\right]\phi = \left[\delta,\pd_\mu\right]
\phi+\left[\delta x^\nu\pd_\nu,\pd_\mu\right]$$
since $\delta$ is a linear operator, it commutes with the 4-derivative; expanding $
\delta x^\nu$ yields
$$\Delta\pd_\mu\phi = \left[\varepsilon^{\nu\rho}x_\rho\pd_\nu,\pd_\mu\right]=
\varepsilon^{\nu\rho}x_\rho\pd_\nu(\pd_\mu\phi) - \varepsilon^{\nu\rho}\pd_\mu(
x_\rho\pd_\nu\phi)$$
expanding with the product rule
$$\Delta\pd_\mu\phi = \varepsilon^{\nu\rho}x_\rho\pd_\nu\pd_\mu\phi - 
\varepsilon^{\nu\rho}x_\rho\pd_\mu\pd_\nu\phi - \varepsilon^{\nu\rho}(\pd_\mu x_\rho)
\pd_\nu\phi$$
the first two terms are zero, because $\varepsilon^{\nu\rho}\left(x_\rho\pd_\nu\pd_\mu
- x_\rho\pd_\mu\pd_\nu\right)=0$, and $\pd_\mu x_\rho = g_{\mu\rho}$, so
$$\Delta\pd_\mu\phi = -\varepsilon^{\nu\rho}g_{\mu\rho}\pd_\nu\phi =
\varepsilon^{\rho\nu}g_{\mu\rho}\pd_\nu\phi.$$
So finally
\begin{align*}
\delta(\pd_\mu\phi) + \varepsilon^{\rho\sigma}x_\rho\pd_\sigma\pd_\mu\phi = 
\varepsilon\indices{^\rho^\nu}g_{\rho\mu}\pd_\nu\phi\\
\delta(\pd_\mu\phi) = \varepsilon\indices{^\rho^\sigma}g_{\mu\rho}
g\indices{^\nu_\sigma}(\pd_\nu\phi)-\varepsilon^{\rho\sigma}x_\rho\pd_\sigma
(\pd_\mu\phi)
\end{align*}
antisymmetrizing both addends with respect to the indices of $\varepsilon$ yields
$$\delta(\pd_\mu\phi) = \frac{1}{2}\varepsilon^{\rho\sigma}\left[
\underbrace{\left(g_{\mu\rho}g\indices{^\nu_\sigma} - g_{\mu\sigma}
g\indices{^\nu_\rho}\right)}_\text{Independent of $x^\mu$}
(\pd_\nu\phi) - \underbrace{\left(x_\rho\pd_\sigma - x_\sigma\pd_\rho\right)}_\text{
$x^\mu$-dependent}(\pd_\mu\phi)\right].$$
This form is convenient, because it is clear that the transformation of the field has
both a $x^\mu$-dependent generator, which is proportional to the generalized angular
momentum $L_{\rho\sigma}$, and a generator that does not depend from $x^\mu$,
which corresponds to a spin operator $\left(S_{\rho\sigma}\right)\indices{_\mu^\nu}$:
this spin operator is none other than the generator of the finite-dimensional Lorentz
transformations $\left(M_{\rho\sigma}\right)\indices{_\mu^\nu}$ and the orbital
angular momentum operator that of the infinite-dimensional transformations.

$$\delta(\pd_\mu\phi) = -\frac{i}{2}\varepsilon^{\rho\sigma}\left(S_{\rho\sigma}
\right)\indices{_\mu^\nu}(\pd_\nu\phi) -
\frac{i}{2}\varepsilon^{\rho\sigma}L_{\rho\sigma}(\pd_\mu\phi)$$

This example underlines the fact that vector fields can also present internal degrees
of freedom, identified as spin, as they are $(n,m)\neq(0,0)$ irreducible
representations of the Poincaré group.
